\subsection{Mapping cylinders around a definable zero locus}
\label{subsec:zero-locus-mapping-cylinders}
Mapping-cylinder neighborhoods of closed semialgebraic sets over real
closed fields are classical.  Delfs and Knebusch prove their existence
in locally complete ambient spaces
\cite[Theorem~2, p.~48; \S3, pp.~63--65]{DelfsKnebusch1984Retractions};
their Proposition~4.1, p.~66, also gives relative homotopy uniqueness
of extensions.  The construction below records the compact definable
version from scalar triangulation, with the specified outer face,
the continuation of a prescribed strict tube, and the empty cases
made explicit.

Let $K$ be a real closed field with a fixed o-minimal expansion.
Compactness in this subsection means definable compactness, equivalently
closedness and boundedness for the affine definable sets used here.
For a continuous definable map $\varphi:B\to Y$, put
\[
 M_\varphi=(B\times[0,1]\amalg Y)/((b,0)\sim\varphi(b)).
\]
An \emph{admissible mapping-cylinder description} of a pair $(V,B)$
with core $Y$ identifies $M_\varphi$ with $V$, is the identity on $Y$,
and identifies $(b,1)$ with the specified point $b$ of the outer face.
The last convention fixes the domain of the attaching map.

\begin{theorem}[Mapping-cylinder descriptions of strict tubes]
\label{thm:zero-locus-mapping-cylinder}
Let $X\subset K^n$ be closed and bounded, and let
$f:X\to K_{\ge0}$ be continuous and definable.  Put $Y=f^{-1}(0)$.
For all sufficiently small $\epsilon>0$, the sublevel set
$V_\epsilon=f^{-1}([0,\epsilon])$ admits an admissible
mapping-cylinder description with outer face
$B_\epsilon=f^{-1}(\epsilon)$.  It is a neighborhood of $Y$ in
$X$, and the displayed product structure off $Y$ extends across
$B_\epsilon$ to an exterior collar.  Its pointed quotient by $Y$
is a rigid cone.

More generally, let $V\subset X$ be a closed strict tube about $Y$ whose
pointed quotient by $Y$ has a specified definable rigid-cone
description with base $B$.  Then $(V,B)$ admits an admissible
mapping-cylinder description, possibly different from the initial
cone description.  For a fixed $(V,B,Y)$, all attaching maps
$B\to Y$ obtained from admissible descriptions are definably homotopic.
\end{theorem}
In the second assertion, the quotient-cone data give a continuous
definable radius $\rho:V\to[0,1]$, with $\rho^{-1}(0)=Y$, and a
specified definable homeomorphism
\[
 e:B\times(0,1]\longrightarrow V\setminus Y,
 \qquad \rho(e(b,t))=t,\qquad e(b,1)=b.
\]
The strict tube also includes an exterior collar of its specified
outer face. Since $V$ is closed in $X$, it is closed and bounded.
These are the data used in the construction below.
\begin{proof}
The triangulation theorem for a continuous definable scalar
function on a compact definable set
\cite[Theorem~2.2(1), p.~6]{Shiota2010} gives a finite polyhedron
$P$, a definable homeomorphism $h:P\to X$, and a nonnegative
piecewise affine function $g=f\circ h$.  Choose a triangulation
on which $g$ is affine on every simplex.  Its zero locus $Z$ is
the subcomplex consisting of faces spanned by zero-valued vertices.
If any positive vertex values occur, choose $\epsilon>0$ with
$2\epsilon$ smaller than all of them.

On a simplex having both zero and positive vertices, write the
barycentric coordinates of $x$ as $(\lambda_v)_v$, and put
\[
 a(x)=\sum_{g(v)=0}\lambda_v,
 \qquad
 y(x)=\frac{1}{a(x)}\sum_{g(v)=0}\lambda_vv.
\]
Whenever $0<g(x)\le2\epsilon$, both the zero and positive
coordinate sums are nonzero.  Indeed
\[
 g(x)=\sum_{g(v)>0}\lambda_vg(v)
   >2\epsilon\sum_{g(v)>0}\lambda_v,
\]
so the latter sum is less than one.  Simplices with only positive
vertices do not meet this sublevel set.  Thus $y(x)$ is defined
there and is continuous across simplex faces: coordinates of
vertices outside a common face are zero, leaving the same formula.

Let $L=g^{-1}(\epsilon)$ and set $\varphi(b)=y(b)$ for $b\in L$.
For $0\le t\le1$ define
\[
 \Psi(b,t)=\varphi(b)+t\bigl(b-\varphi(b)\bigr).
\]
This lies in the simplex containing $b$ and satisfies
$g(\Psi(b,t))=t\epsilon$.  Together with the identity on $Z$,
it defines a continuous definable map $M_\varphi\to g^{-1}([0,\epsilon])$.
For a point $x$ with $r=g(x)>0$, the inverse follows the same
segment to height $\epsilon$:
\[
 b=y(x)+\frac{\epsilon}{r}\bigl(x-y(x)\bigr),
 \qquad t=r/\epsilon.
\]
The total positive barycentric weight of $b$ is
$\epsilon/g(z)<1$, where $z$ is the normalized positive-face
component of $x$; hence $b$ remains in that simplex.  The normalized
zero-face component stays equal to $y(x)$, so these formulas are
mutual inverses off $Z$.  On $Z$ the mapping-cylinder equivalence
relation gives exactly the identity.  A continuous definable
bijection from a closed bounded definable model of this mapping
cylinder to the Hausdorff set $g^{-1}([0,\epsilon])$ is a
homeomorphism.  For example, a closed bounded model is supplied
by $(tb,\varphi(b),t)$ together with $(0,y,0)$ for $y\in Z$.

The same segments continue to height $2\epsilon$: the positive
vertex values are larger than that height, so the same inverse
formula applies.  Consequently
\[
 L\times(0,2\epsilon]\longrightarrow g^{-1}((0,2\epsilon]),
 \qquad (b,r)\longmapsto
 \varphi(b)+(r/\epsilon)(b-\varphi(b))
\]
is a definable product description.  It provides both the inner
collar and the exterior collar at height $\epsilon$.
The set $g^{-1}([0,\epsilon))$ is open in $P$ and contains $Z$,
so this is a neighborhood of the core.  Transport by $h$ proves
the first assertions.  Collapsing the core of a mapping cylinder
gives the cone on its outer face.

The exceptional cases are included.  Components contained in $Z$
need no outer-face preimage.  If $Y=X$, take $B_\epsilon=\varnothing$
and the mapping cylinder $Y$.  If $Y=\varnothing$, then either
$X=\varnothing$ or $f$ has a positive minimum; choose $\epsilon$
below that minimum and the neighborhood and outer face are empty.
For quotients use the pointed convention $X/\varnothing=X\amalg\{o\}$,
so the empty-neighborhood quotient is the cone on the empty set.

For the prescribed strict tube, its cone radius pulls back to a
continuous definable function $\rho:V\to[0,1]$ with zero set $Y$.
The given rigid cone identifies $V\setminus Y$ with
$B\times(0,1]$, preserving $\rho$.  Apply the construction just
proved to obtain an admissible mapping cylinder below a small
height $\eta>0$.  Let $j:B\to\rho^{-1}(\eta)$ be the restriction
of the given positive-radius product.  Transport the new attaching
map along $j$.  Below height $\eta$, use the new mapping cylinder
with radial coordinate rescaled by $\eta$; above height $\eta$,
use the original product.  Both formulas equal $j(b)$ on the
common level.  They glue to a continuous definable bijection
$M_\varphi\to V$ that is the identity on the core and outer face,
and definable compactness again makes it a homeomorphism.
This proves the assertion without requiring the original cone
rays to extend continuously to individual points of $Y$.

Finally, two admissible descriptions give strong definable
deformation retractions $H_i:V\times[0,1]\to V$ onto $Y$,
with final retractions $r_i:V\to Y$.  The homotopy
\[
 (x,s)\longmapsto r_0(H_1(x,s))
\]
runs from $r_0$ to $r_1$ and fixes $Y$ pointwise, since
$r_0|_Y$ is the identity.  Its restriction to the fixed outer
face $B$ joins the two attaching maps.  This proves independence
of their homotopy class; it makes no assertion that the space of
descriptions is contractible.
\end{proof}
